\documentclass[11pt]{article}
\usepackage[margin=1in]{geometry}
\usepackage{amsmath,amssymb,amsthm,mathtools}
\usepackage{xcolor}\usepackage[colorlinks=true,linkcolor=violet]{hyperref}
\newtheorem{lemma}{Lemma}\newtheorem{proposition}{Proposition}\newtheorem{theorem}{Theorem}
\theoremstyle{remark}\newtheorem{remark}{Remark}
\newcommand{\Z}{\mathbb{Z}}\newcommand{\Rr}{\mathbb{R}}\newcommand{\var}{\operatorname{var}}
\newcommand{\Lop}{\mathcal L}\newcommand{\e}{\mathbf e}
\title{\bfseries UNI verification for the log-growth roof, and the joint mixed-LLT operator}
\author{BB(6) research program}\date{2026-09-08}
\begin{document}\maketitle
\begin{abstract}
On the suspension of the induced Gauss-type map $G$ by the log-growth roof
$\tau$, we (i) verify the uniform-non-integrability (UNI) condition for $\tau$
--- rigorously, from an exact non-proportionality of the branch data --- and
(ii) set up the joint transfer operator whose local limit theorem yields the
native-resolution equidistribution that the density form requires. UNI is the
hypothesis Dolgopyat's theorem needs; with it verified, the uniform spectral
gap for the additive roof follows from settled theory. We are explicit that a
tempting shortcut (cohomology of $\tau$ to the geometric potential) is
\emph{false}, and that the uniform gap and the mixed LLT are then invoked, not
re-proved.
\end{abstract}

\section{The model roof}
In the scale-invariant model, a section state is $(a,b)$ with $\xi=b/a$ and
offset $o=b-a=a(\xi-1)$. A return in branch $k$ applies the integer matrix
$P_k=\begin{psmallmatrix}p_k&q_k\\ r_k&s_k\end{psmallmatrix}=C_\ell^{\,k}F_\ell$,
so $a'=p_ka+q_kb$, $b'=r_ka+s_kb$, and $o'=b'-a'=(r_k-p_k)a+(s_k-q_k)b
=a\,L_k(\xi)$ with
\[
L_k(\xi):=(r_k-p_k)+(s_k-q_k)\,\xi .
\]
Hence the log-growth roof is, exactly,
\begin{equation}\label{eq:tau}
\tau(\xi,k)=\log\Big|\tfrac{o'}{o}\Big|=\log|L_k(\xi)|-\log|\xi-1|.
\end{equation}
The term $-\log|\xi-1|$ is branch-independent, a function of the source alone;
it is a continuous coboundary and drops out of every temporal-distance
comparison.

\begin{remark}[a shortcut that fails]
One might hope $\tau$ is cohomologous to $-\tfrac12\log|G'|+\text{(locally
constant)}$, reducing UNI to the geometric potential. It is \emph{not}: the
identity $\tau\overset?=-\tfrac12\log|G'|+\tfrac12(k{+}1)\log3+(\log z'-\log
z)$ fails numerically (mean discrepancy $1.32$, correlation $0.57$ over
$1.3\times10^6$ returns). $\tau$ must be treated as its own roof.
\end{remark}

\section{UNI, verified}
\begin{proposition}[Non-degenerate temporal distance]\label{p:uni}
For every pair of distinct branches $i\neq j$ the temporal-distance function
\[
\Delta_{ij}(\xi):=\big(\tau(\xi,i)-\tau(\xi,j)\big)=\log|L_i(\xi)|-\log|L_j(\xi)|
\]
is non-constant; equivalently $\tau$ is not cohomologous to a function
constant on branches. Consequently $\tau$ satisfies the uniform
non-integrability condition.
\end{proposition}
\begin{proof}
By \eqref{eq:tau} the common $-\log|\xi-1|$ cancels, leaving
$\Delta_{ij}=\log|L_i|-\log|L_j|$. The bottom rows of $P_k=C_\ell^kF_\ell$ give
$(r_k-p_k,\,s_k-q_k)=(-4,1),(7,-4),(-19,7),(40,-19),(-97,40),(217,-97),\dots$,
whose slope ratios $(s_k-q_k)/(r_k-p_k)=-\tfrac14,-\tfrac47,-\tfrac{7}{19},
-\tfrac{19}{40},\dots$ (the $-\rho_k$) are pairwise distinct. Thus $L_i,L_j$
are non-proportional affine forms, so $\Delta_{ij}'(\xi)=\frac{s_i-q_i}{L_i(\xi)}
-\frac{s_j-q_j}{L_j(\xi)}\not\equiv0$ (a nonzero rational function; e.g.
$\Delta_{02}'$ takes $0.9,-4.5,0.5625$ at $\xi=2,3,5$). Non-constancy of every
$\Delta_{ij}$ is exactly the statement that no single continuous $\varphi$ and
constants $c_k$ solve $\tau(\cdot,k)=\varphi\circ G_k-\varphi+c_k$; since the
branches are real-analytic (M\"obius), non-cohomology to locally constant is
the UNI condition in the sense of Dolgopyat/Naud/Stoyanov.
\end{proof}

\begin{proposition}[Uniform gap, from UNI]\label{p:gap}
$G$ is a real-analytic full-branch expanding Markov map with a spectral gap,
and $\tau$ is a real-analytic roof satisfying UNI (Prop.~\ref{p:uni}) and
non-lattice (its increment law has characteristic function bounded away from
$1$ off the origin, verified). By Dolgopyat's theorem for such systems
\emph{(cited: Dolgopyat 1998; Naud 2005; Stoyanov 2011)}, the twisted operator
$\Lop_{i\xi}f=\sum_k|g_k'|^{s^\ast}\e(\xi\tau_k)f(g_k\cdot)$ has spectral radius
bounded away from $1$ uniformly for $|\xi|\ge\varepsilon$:
$\sup_{|\xi|\ge\varepsilon}\|\Lop_{i\xi}^{\,n}\|\le C\theta_\varepsilon^{\,n}$,
$\theta_\varepsilon<1$.
\end{proposition}
This is the additive-roof gap that was unavailable for the multiplicative
cost; it is invoked here, its hypotheses (UNI, non-lattice, analytic gap at
$\xi=0$) now discharged.

\section{The joint operator and the local limit theorem}
Let $\Lop_{i\xi}$ act on $BV[0,1]$ as above. Its leading eigenvalue $\lambda(i\xi)$
defines the pressure $P(i\xi)=\log\lambda(i\xi)$, analytic near $0$ with
\[
P(0)=0,\qquad P'(0)=i\Lambda,\qquad P''(0)=-\sigma^2,\quad \Lambda\approx0.51,\ \sigma^2\approx0.265>0.
\]
The pair $(z_d,\,\ell_d)$, $\ell_d=\log\text{scale}_d=\sum_{j<d}\tau(z_j)$,
is governed by $\Lop_{i\xi}$: for $\varphi\in BV$,
$\mathbb E_z[\e(\xi\ell_d)\,\psi(z_d)]=\langle\psi,\Lop_{i\xi}^{\,d}\varphi\rangle$.
\begin{theorem}[Mixed LLT; cited assembly]\label{t:llt}
Under Prop.~\ref{p:gap} and $\sigma^2>0$, the Nagaev--Guivarc'h method
(Aaronson--Denker for Gibbs--Markov maps) gives a local limit theorem: for
bounded intervals $J$ and $\psi\in BV$,
\[
\sqrt{d}\;\mathbb E_z\big[\mathbf 1_{\{\ell_d-\Lambda d\in J\}}\,\psi(z_d)\big]
\;\longrightarrow\;\frac{|J|}{\sqrt{2\pi\sigma^2}}\int\psi\,d\mu,
\]
uniformly in the centering, and jointly with the equidistribution of $z_d$
w.r.t.\ $\mu$. In particular $\ell_d\bmod\log M$ equidistributes, jointly with
$z_d$, at a polynomial-in-$M$ rate.
\end{theorem}

\section{From the joint LLT to native-resolution equidistribution}
At a crossing of scale $S$, the offset is $o_d=z_d\,e^{\ell_d}$; the fatal
event $o_d\in\{1,3,5\}$ is, to resolution $M\sim S$, a joint condition on
$(z_d\bmod\text{-},\ell_d\bmod\log M)$. Theorem~\ref{t:llt} supplies exactly
the joint equidistribution needed: the per-crossing hazard is
$6h(0)/S+o(1/S)$ with the error controlled by the LLT rate, uniformly in $S$.
This is the content of the open uniform-gap lemma $(\star)$, now delivered
through the additive-roof route.

\section{Status}
\emph{Rigorous here:} the model roof \eqref{eq:tau}; the UNI condition for
$\tau$ via the exact non-proportionality of the branch forms
(Prop.~\ref{p:uni}); the pressure data $P(0)=0$, $\sigma^2>0$; the falsity of
the geometric-potential shortcut. \emph{Invoked (settled theory, hypotheses
discharged):} the uniform gap from UNI (Dolgopyat/Naud/Stoyanov,
Prop.~\ref{p:gap}); the mixed LLT (Nagaev--Guivarc'h / Aaronson--Denker,
Thm.~\ref{t:llt}). \emph{Still to write:} the $O(1/S)$ model-to-exact
corrections carried through the LLT (via the truncation of the formal
supplement); and the fully detailed Nagaev--Guivarc'h perturbation for the
joint operator.

The open lemma is therefore no longer ``outside the theory needing a new
idea.'' Its one non-standard hypothesis --- UNI for a genuinely
multiplicative-looking cost --- has been reduced, via the suspension, to UNI
for an additive real-analytic roof, and \emph{that} is now verified from the
exact branch arithmetic. What remains is the invocation of two standard
theorems whose hypotheses are discharged, plus routine perturbation
bookkeeping. No closure is claimed, but the analytic obstruction is removed.
\end{document}
